\section{Related Work}

Plant disease detection has evolved from traditional image processing and handcrafted feature-based methods toward deep learning-based approaches. Earlier computer vision techniques typically relied on preprocessing, segmentation, color or texture feature extraction, and machine learning classifiers. Barbedo~\cite{barbedo2013} reviewed digital image processing methods for detecting and classifying plant diseases, while Zhang et al.~\cite{zhang2014} discussed the use of computer vision in plant pathology. Although these methods demonstrated the feasibility of automated disease diagnosis, their reliance on handcrafted features and controlled imaging conditions limited their robustness across different crops, disease symptoms, and environmental variations. With the availability of plant disease image datasets, convolutional neural networks (CNNs) became widely adopted for automated feature learning. Mohanty et al.~\cite{mohanty2016} demonstrated the effectiveness of deep learning on the PlantVillage dataset, while Sladojevic et al.~\cite{sladojevic2016} and Ferentinos~\cite{ferentinos2018} further showed that CNN-based models can achieve strong performance in plant disease classification. However, many high-performing CNNs are computationally demanding, which creates challenges for mobile and resource-constrained deployment.

Recent research has increasingly focused on lightweight CNN architectures to reduce computational complexity while maintaining competitive classification performance. MobileNet introduced depthwise separable convolutions as an efficient alternative to standard convolutional operations~\cite{howard2017mobilenets}, and MobileNetV2 further improved this design using inverted residual blocks and linear bottlenecks~\cite{sandler2018mobilenetv2}. Other architectures, such as InceptionV3~\cite{szegedy2016rethinking}, Xception~\cite{chollet2017xception}, and EfficientNet~\cite{tan2019efficientnet}, have also contributed to efficient CNN design, although they differ significantly in model size, parameter count, and deployment suitability. In plant disease detection, practical systems must consider not only accuracy but also inference time, memory requirements, and model size. Khan et al.~\cite{khan2023edge} investigated plant disease detection for edge computing devices and emphasized the importance of optimization and quantization for resource-constrained platforms. Similarly, Duhan et al.~\cite{duhan2025rtr} proposed RTR\_Lite\_MobileNetV2, an enhanced lightweight MobileNetV2-based model for plant disease detection and classification, showing the growing interest in MobileNetV2 variants for efficient agricultural diagnosis. Li et al.~\cite{li2024tea_mobilenetv2} also improved MobileNetV2 for tea leaf disease identification using Coordinate Attention, multi-branch convolution, and model compression, further supporting the relevance of lightweight CNNs for field-oriented applications.

Attention mechanisms have become an important strategy for improving CNN feature representation without necessarily replacing lightweight backbones with heavier architectures. The squeeze-and-excitation network introduced channel-wise feature recalibration~\cite{hu2018senet}, while the Convolutional Block Attention Module (CBAM) extended this idea by sequentially applying channel and spatial attention~\cite{woo2018cbam}. This is particularly relevant for plant disease recognition because symptoms often appear as localized spots, discoloration, or texture irregularities on leaf surfaces. Recent attention-based studies have further shown the usefulness of attention mechanisms in agricultural image analysis. Subburaj et al.~\cite{subburaj2025plantention} proposed Plantention, a lightweight attention-based model using a MobileNetV2 encoder for multi-crop leaf disease classification. Wen et al.~\cite{wen2025mnasnet} introduced a lightweight CNN for tea leaf disease and pest recognition and compared several attention mechanisms, including CA, ECA, SE, and CBAM. Pal et al.~\cite{pal2025mobres} proposed a lightweight and explainable CNN model for plant disease diagnosis and used interpretability techniques such as Grad-CAM, Grad-CAM++, and LIME to analyze model predictions. These studies indicate that attention-enhanced lightweight CNNs are becoming increasingly important for accurate, interpretable, and deployment-aware plant disease detection.

Despite these advances, many plant disease detection studies still emphasize classification accuracy more strongly than deployment-oriented evaluation. Recent review studies note that practical plant disease recognition systems must address challenges such as variable illumination, symptom similarity, overfitting, dataset bias, limited generalization, and real-world deployment constraints~\cite{zhao2025review,nyawose2025review,sajitha2024review}. In particular, fewer studies jointly analyze classification metrics such as accuracy, precision, recall, F1-score, and AUC together with computational indicators such as parameter count, model size, TensorFlow Lite size, and inference latency. This creates a gap between high-performing laboratory models and practical mobile or edge-based agricultural diagnostic systems. To examine this trade-off, this paper evaluates CBAM after the final MobileNetV2 feature map and before global average pooling. The objective is to improve the baseline while retaining a compact candidate architecture; suitability for a particular resource-constrained device is not asserted without device-specific validation.

Table~\ref{tab:recent_lightweight_work} positions the present study against representative work published in 2024--2025. The comparison is qualitative because the studies use different datasets, splits, preprocessing, hardware, and reporting conventions; their accuracy and latency values should not be treated as a controlled leaderboard. Importantly, recent work already combines MobileNetV2 with attention or other lightweight refinements. The distinction of the present study is therefore the evaluated single late-CBAM placement and its measured incremental trade-off against an unmodified MobileNetV2 baseline, not the invention of an attention-enhanced MobileNet family.

\begin{table}[H]
\caption{Scope comparison with representative recent lightweight plant-disease studies.}
\label{tab:recent_lightweight_work}
\centering
\small
\begin{tabularx}{\textwidth}{p{2.8cm} p{4.3cm} X}
\toprule
\textbf{Study} & \textbf{Lightweight strategy} & \textbf{Scope and relation} \\
\midrule
Li et al. (2024)~\cite{li2024tea_mobilenetv2} & MobileNetV2, coordinate attention, multi-branch convolution, and compression & Tea-leaf evaluation; a broader modification using a different attention mechanism. \\
Duhan et al. (2025)~\cite{duhan2025rtr} & RTR\_Lite\_MobileNetV2 & Confirms that lightweight MobileNetV2 variants form an active, non-unique design family. \\
Subburaj et al. (2025)~\cite{subburaj2025plantention} & Attention model with a MobileNetV2 encoder & Multi-crop scope; this study instead isolates one late CBAM block. \\
Wen et al. (2025)~\cite{wen2025mnasnet} & Lightweight CNN comparing CA, ECA, SE, and CBAM & Provides an attention ablation that remains necessary here under one common protocol. \\
Pal et al. (2025)~\cite{pal2025mobres} & MobileNetV2 plus residual learning (3.51 M parameters) & Two datasets and Grad-CAM/Grad-CAM++/LIME; 99.47\% on PlantVillage, under a different protocol. \\
This study & One CBAM block after the final MobileNetV2 convolutional map & Reports an incremental accuracy--cost trade-off; device and interpretability evidence remain outstanding. \\
\bottomrule
\end{tabularx}
\normalsize
\end{table}
